Ultra-long-context capability is becoming indispensable for frontier LLMs: agentic workflows, repository-scale code reasoning, and persistent memory all require the model to jointly attend over hundreds of thousands to millions of tokens—yet the quadratic cost of softmax attention makes this untenable at deployment scale. We introduce MiniMax Sparse Attention (\method{}), a blockwise sparse attention built upon Grouped Query Attention (GQA).
A lightweight Index Branch scores key–value blocks and independently selects a Top-$k$ subset for each GQA group, enabling group-specific sparse retrieval while maintaining efficient block-level execution; the Main Branch then performs exact block-sparse attention over only the selected blocks. Designed around a principle of simplicity and scalability, \method{} is deliberately streamlined, making it straightforward to deploy efficiently across a broad range of GPUs.
To translate sparsity into practical speedups, we co-design \method{} with a GPU execution path that uses exp-free Top-$k$ selection and KV-outer sparse attention to improve tensor-core utilization under block-granular access. On a 109B-parameter model with native multimodal training, \method{} performs on par with GQA while reducing per-token attention compute by $28.4\times$ at 1M context. Paired with our co-designed kernel, \method{} achieves $14.2\times$ prefill and $7.6\times$ decoding wall-clock speedups on H800. Our inference kernel is available at: \url{https://github.com/MiniMax-AI/MSA}. A production-grade natively multimodal model powered by \method{} has been publicly released at: \url{https://huggingface.co/MiniMaxAI/MiniMax-M3}.  

